\documentclass[11pt]{article}
\usepackage[margin=0.7in]{geometry}
\usepackage{amsmath,amssymb}
\usepackage{enumitem}
\usepackage{tikz}
\usepackage{pgfplots}
\pgfplotsset{compat=1.18}
\usepackage{fancyhdr}
\usepackage{parskip}
\setlist[enumerate]{leftmargin=*, itemsep=6pt}
\pagestyle{fancy}
\fancyhf{}
\fancyfoot[L]{\small\textit{en}Vision\textsuperscript{\textregistered} Algebra 2}
\fancyfoot[C]{\small 5-1 Lesson Quiz, short form}
\fancyfoot[R]{\small\thepage}
\renewcommand{\headrulewidth}{0pt}
\newcommand{\choice}[1]{\tikz[baseline=-0.6ex]\node[draw, circle, inner sep=0.6pt, minimum size=3.3mm, font=\scriptsize\bfseries] {#1};\;}
\newcommand{\checkbox}{$\square$\;}
\pgfplotsset{quizgrid/.style={
  axis lines=middle, grid=both, grid style={gray!45},
  ticklabel style={font=\scriptsize, fill=white, inner sep=1pt},
  xlabel={$x$}, ylabel={$y$}, axis line style={thick}, clip=true}}
\renewcommand{\arraystretch}{1.6}
% one "circle one" menu for the TE's drop-down choices in item 3
\newcommand{\menu}[1]{\fbox{\small #1}}

\begin{document}

\noindent\textbf{Name}\; \rule{6cm}{0.4pt} \hfill Period \rule{1.2cm}{0.4pt} \quad Date \rule{2.2cm}{0.4pt}

\vspace{4pt}\noindent\rule{\textwidth}{0.6pt}

\vspace{2pt}\noindent{\Large\textbf{5-1 \; Lesson Quiz: nth Roots, Radicals, and Rational Exponents}} \hfill {\small short form}

\vspace{4pt}\noindent\rule{\textwidth}{0.6pt}

{\small Question numbers match the Savvas Lesson Quiz. Questions 1 and 4 are not on this form. Show your work.}

\vspace{4pt}
\begin{minipage}[t]{0.52\textwidth}
\textbf{2.} Select \textbf{all} expressions that are equivalent to $x^{\frac34}$.

\vspace{4pt}
\begin{itemize}[leftmargin=1.2em, itemsep=7pt]
\item[] \checkbox \textbf{A.} $\left(x^3\right)^{\frac14}$
\item[] \checkbox \textbf{B.} $x^{\frac86}$
\item[] \checkbox \textbf{C.} $\left(x^{\frac14}\right)^{\frac24}$
\item[] \checkbox \textbf{D.} $\left(x^{\frac14}\right)^3$
\item[] \checkbox \textbf{E.} $\sqrt[3]{x^4}$
\item[] \checkbox \textbf{F.} $\sqrt[4]{x^3}$
\end{itemize}
\end{minipage}\hfill
\begin{minipage}[t]{0.44\textwidth}
\textbf{Work:}

\vspace{4pt}\rule{\linewidth}{0.4pt}

\vspace{14pt}\rule{\linewidth}{0.4pt}

\vspace{14pt}\rule{\linewidth}{0.4pt}

\vspace{14pt}\rule{\linewidth}{0.4pt}
\end{minipage}

\vspace{10pt}\noindent\rule{\textwidth}{0.4pt}

\vspace{6pt}
\textbf{3.} Explain what $243^{\frac35}$ means, then evaluate it. \ {\small(Circle one word in each box.)}

\vspace{6pt}
$243^{\frac35}$ is the \menu{square \;/\; cube \;/\; fourth \;/\; fifth} root of 243

\vspace{4pt}\hspace*{2em}\menu{squared \;/\; cubed \;/\; to the fourth power \;/\; to the fifth power}.

\vspace{6pt}
Its value is \menu{3 \;/\; 9 \;/\; 27 \;/\; 81}.

\vspace{4pt}
\textbf{Work:}

\vspace{4pt}\rule{\linewidth}{0.4pt}

\vspace{14pt}\rule{\linewidth}{0.4pt}

\vspace{10pt}\noindent\rule{\textwidth}{0.4pt}

\vspace{6pt}
\begin{minipage}[t]{0.52\textwidth}
\textbf{5.} Solve the equation $-4x^3=32$.

\vspace{10pt}
$x=$ \rule{3cm}{0.4pt}
\end{minipage}\hfill
\begin{minipage}[t]{0.44\textwidth}
\textbf{Work:}

\vspace{4pt}\rule{\linewidth}{0.4pt}

\vspace{14pt}\rule{\linewidth}{0.4pt}

\vspace{14pt}\rule{\linewidth}{0.4pt}

\vspace{14pt}\rule{\linewidth}{0.4pt}
\end{minipage}

\vfill
\noindent\textbf{Scoring:} Question 3 and Question 5 are 10 points each. Question 2 is a bonus of up to 5 points (all three correct boxes, none wrong). Total out of 20.

\newpage
\noindent{\LARGE\textbf{5-1 Lesson Quiz, short form: answer key}}

\vspace{4pt}\noindent\rule{\textwidth}{0.6pt}

\begin{enumerate}
\item[\textbf{2.}] \textbf{A, D, F.} A: $\left(x^3\right)^{\frac14}=x^{\frac34}$. D: $\left(x^{\frac14}\right)^3=x^{\frac34}$. F: $\sqrt[4]{x^3}=x^{\frac34}$ (index 4 is the denominator, power 3 the numerator). Not B: $x^{\frac86}=x^{\frac43}$. Not C: $\left(x^{\frac14}\right)^{\frac24}=x^{\frac18}$. Not E: $\sqrt[3]{x^4}=x^{\frac43}$ (index and power swapped). Savvas quiz item 2, DOK 2, skill A82; taught by Example 2, rehearsed by the packet's Practice. \textbf{Bonus 5:} A, D and F with no extra box; 3 for two of the three and none wrong; 0 otherwise.
\item[\textbf{3.}] \textbf{fifth; cubed; 27.} $243^{\frac35}$ is the fifth root of 243 cubed: $\left(\sqrt[5]{243}\right)^3=3^3=27$, because $3^5=243$. Savvas quiz item 3, DOK 2, skill A82; taught by Example 2 (Try It 2 is the same task). \textbf{10 points:} 3 for \emph{fifth}, 3 for \emph{cubed}, 4 for 27 with work (243 written as $3^5$, or the fifth root found first).
\item[\textbf{5.}] \textbf{$x=-2$.} Divide by $-4$: $x^3=-8$. Take the cube root of both sides (raise to the $\frac13$ power): $x=\sqrt[3]{-8}=-2$. Check: $-4(-2)^3=-4(-8)=32$ \checkmark. The index 3 is odd, so there is exactly one real solution, and it is negative. Savvas quiz item 5, DOK 2, skill A84; taught by Example 5 (Try It 5a). \textbf{10 points:} 4 for $x^3=-8$, 3 for the cube root step, 3 for $x=-2$. Common error: $x^3=-128$ (multiplying instead of dividing) or $x=\pm2$ (the even-index $\pm$ used on an odd index).
\end{enumerate}

\textbf{Items left off this form} (available for a retake or Form B): Savvas item 1, the real fourth roots of 256 (answer 4 and $-4$; DOK 1, A83; Example 1, cut, though the $\pm$ move for an even index is Example 5B); Savvas item 4, simplify $\sqrt[4]{81m^{12}n^4}$ with all variables positive (answer C, $3m^3n$; DOK 1, A80; Example 4, cut in the fast packet).

\textbf{IXL replacement:} SmartScore 70 or above on 9TM, KJX and U42 replaces this quiz with 20/20.

\end{document}
